\subsection{Summary}

We compared eighteen models for 7-day-ahead daily forecasting of EV
charging demand at N~Boulder Rec~1, the busiest public station in
Boulder, Colorado. The leaderboard goes from naive baselines through
the full classical (S)ARIMA family, an LSTM with hyper-parameter
sweep, four zero-shot foundation models (TimesFM 1.0, TimesFM 2.0,
Chronos-T5 base, Chronos-T5 large), and finally a from-scratch
PatchTST transformer with RevIN, channel-mixed exogenous variables
and a three-seed ensemble. The best model on the leaderboard is a
simple-mean ensemble of PatchTST and LSTM, with MAE
$25.09 \pm 9.36$~kWh. The next best is PatchTST alone (MAE
$25.19 \pm 8.90$~kWh, lowest variance among forecasting models),
followed by ARIMAX (2,1,3) at $25.52 \pm 14.38$~kWh.

\subsection{What We Beat and What We Did Not}

The strongest claim we can make honestly is: \emph{within this
project's evaluation protocol --- same series, same train/test split,
same rolling-origin windows, same metric --- the PatchTST + LSTM
ensemble beats the strongest non-transformer model (ARIMAX) on both
the mean MAE and the standard deviation, with the standard deviation
roughly 35\% lower.} That is a real improvement.

The weaker claim we \emph{cannot} make is that we beat the published
state of the art. The three peer-reviewed papers that use this dataset
(Koohfar~2023~\cite{koohfar2023transformer},
Wu/Chowdhury~2025~\cite{wu2025catformer},
Kyriakopoulos~2025~\cite{kyriakopoulos2025comparative}) all report
their numbers in normalised units --- z-score for Koohfar and
Kyriakopoulos, MinMax for Wu --- on different station selections
(aggregated or per-station, but always a different mix from ours) and
mostly different horizons (1 hour, 1 day, or 30/60/90 days, but rarely
7 days). Our 25~kWh and their 0.7--0.9 (normalised) are not in the
same unit system; ``$25 < 0.91$'' is not a meaningful comparison.

The qualitative ordering, on the other hand, matches: transformer
$>$ LSTM $>$ classical, with the gap explained by the exogenous
variables and the long-horizon-friendly architecture. We sit in the
SoTA \emph{family}; we have not run a like-for-like benchmark against
its representatives.

\subsection{Lessons Worth Writing Down}

Three observations from this project that we want to remember:
\begin{enumerate}[leftmargin=*]
    \item \textbf{Seed ensembles work.} A single PatchTST scored
    MAE 27.1~kWh, which is worse than ARIMAX. Three independent
    PatchTSTs (different seeds, otherwise identical) averaged at
    inference scored 25.19~kWh, beating ARIMAX. Almost two kWh of
    variance reduction for the cost of training three models instead
    of one --- the cheapest performance boost in the entire project.
    \item \textbf{Bigger zero-shot $\neq$ better.} Chronos-T5 large
    (710M parameters) is the \emph{worst} foundation model in our
    leaderboard, beaten by its 200M sibling. More parameters means
    the pre-training prior is more strongly amplified at inference,
    which only helps if the prior is correct for the new dataset.
    For Boulder EV demand it is not. TimesFM transfers better
    out-of-the-box.
    \item \textbf{Stacking needs a real validation set.} Our
    leave-one-window-out Ridge and NNLS stacking attempts both failed
    to generalise from 21 days to 7 days. With 28 total test points
    the meta-learner cannot pick up reliable weights. The dumbest
    convex combination --- an unweighted mean of the two strongest
    deep models, picked a priori --- ended up winning. Lesson: do not
    try to learn an ensemble weight without enough data to support
    it.
\end{enumerate}

\subsection{Limitations and Future Work}

The honest limitations of this project are:
\begin{itemize}[leftmargin=*]
    \item \textbf{Single station, 30-day test set.} 28 test points is
    enough to rank models but not enough for confident generalisation
    claims. A future iteration should retrain on each of the top
    stations (Baseline St~1, Carpenter Park~1, \dots) and report a
    panel of results.
    \item \textbf{No direct comparison with published numbers.} As
    discussed above, the published Boulder papers use different units
    and station selections. A defensible ``we beat CAT-Former'' or
    ``we beat Koohfar'' claim would require replicating one of those
    exact setups (same stations, same MinMax / z-score normalisation,
    same horizon) and reporting in their unit system. That is the
    natural next experiment.
    \item \textbf{No probabilistic forecasting.} All our models
    produce point forecasts; we did some basic prediction-interval
    work in an auxiliary notebook (07) but did not push it
    aggressively. For grid planning, calibrated quantiles would be
    much more useful than point forecasts.
    \item \textbf{PatchTST uses calendar-known exogenous only.} We do
    not feed the model a temperature \emph{forecast} for the future
    window; we feed the actual recorded temperature. In production
    the operator would only have a forecast, and adding noise to the
    temperature input would tell us how robust the gain is.
\end{itemize}

The most impactful follow-up would be the first item combined with
the second: re-run PatchTST and the ensemble on the three CAT-Former
stations with MinMax normalisation and the 1-day horizon, and report
in normalised MAE/MSE so the comparison with CAT-Former is honest.
Our PatchTST implementation already supports the necessary changes;
it is essentially a preprocessing job.

\subsection{Closing Note}

This project started as a comparison between SARIMA and LSTM and ended
up as a tour through every major time-series modelling family that
matters in 2026. The verdict is uncomfortable for the loudest hype:
foundation models did not win, and the strongest classical model is
within one kWh of the best deep learning model. The win came from
combining a small purpose-trained transformer (PatchTST) with a
specifically-tuned LSTM, and from running each model three times to
average out initialisation noise. None of those tricks are exotic.
They were, however, the things that moved the needle.
